% ============================================================
%  Section 5 -- Discussion
% ============================================================

\chapter{Discussion}

\section{Interpretation du resultat principal}

Le resultat principal ne doit pas etre lu comme une victoire definitive du
systeme multi-agents, mais comme une indication experimentale forte. Sur les
10 tickets evalues, le SMA produit un meilleur score qualite, consomme beaucoup
moins de tokens, coute moins cher, repond plus vite et obtient un score
composite nettement superieur. Ces gains sont coherents avec l'intuition : un
agent qui interroge directement une base de connaissances ou un statut serveur
peut eviter de longues reponses generatives et fonder sa reponse sur un contexte
plus cible.

Cependant, le score qualite moyen de 3,8/10 reste faible. Le gate de qualite a
donc joue son role : il empeche de transformer une amelioration relative en
decision de deploiement. Dans un contexte industriel, cette rigueur est saine.
Un agent moins mauvais que la reference n'est pas automatiquement exploitable par
un support IT.

\section{Pourquoi la qualite reste faible}

Plusieurs hypotheses expliquent le score qualite observe.

\subsection{Modele local limite}

La campagne principale utilise \code{phi3:latest}. Ce choix est interessant pour
l'execution locale, mais il peut limiter la capacite de synthese et de suivi
d'instructions par rapport a des modeles plus puissants. Les scores historiques
montrent que \code{gpt-3.5-turbo} atteint 7,01/10 sur un autre protocole, ce qui
suggere que le modele est un facteur majeur.

\subsection{Corpus et references incomplets}

Une architecture RAG depend fortement de son corpus. Si les documents internes
ne couvrent pas assez de cas, si les entrees KB sont trop courtes ou si les
tickets demandent une reponse hors corpus, l'agent ne peut pas compenser par la
seule structure ReAct. Les resultats MCP rapides mais moins bons en pertinence
dans \file{benchmark_mcp_tools.json} vont dans ce sens : l'outil fonctionne, mais
la pertinence mesuree reste perfectible.

\subsection{Prompts et format de sortie}

L'agent produit une reponse finale en francais, mais le prompt pourrait etre
renforce pour imposer une structure plus proche des criteres du juge :
diagnostic, cause probable, actions par priorite, verification, escalade et
limites. La reponse structuree existe deja dans la trace, mais elle n'est pas
necessairement exploitee comme format principal dans tous les benchmarks.

\section{Apport de la politique composite}

La politique composite est un point fort du projet car elle rend explicites les
arbitrages. Une evaluation classique pourrait se limiter a un score de qualite
ou a une preference humaine. \bibopsname{} ajoute des dimensions qui comptent
pour l'exploitation : cout, latence, securite et GreenOps.

Le poids de 35 \% donne a la securite est volontairement eleve. Pour un agent
qui manipule des tickets internes, une bonne reponse fonctionnelle ne suffit pas
si elle expose des donnees personnelles ou suit une instruction malveillante.
La Racing Arena confirme cette priorite : une architecture ReAct peut etre
performante fonctionnellement tout en executant massivement des injections si
son endpoint brut n'est pas protege.

La limite de cette politique est la calibration. Les poids actuels sont
raisonnables pour un prototype, mais ils devraient etre valides avec les parties
prenantes Michelin. Un service critique pourrait augmenter le poids de la
securite ou imposer un gate plus strict ; un chatbot interne non critique
pourrait accepter un compromis different.

\section{Robustesse et securite agentique}

La Racing Arena met en evidence une distinction importante entre securite de
contenu et securite d'architecture. Les detecteurs de PII, secrets ou injection
sont necessaires, mais ils ne suffisent pas si l'agent expose des routes faibles,
relaie des prompts sans filtrage ou confond messages de donnees et instructions.

Team B illustre ce risque : l'equipe dispose d'un comportement ReAct, mais son
serveur \code{/query} appelle directement le LLM sur un payload externe. Dans la
course observee, elle execute 93 \% des injections recues. Team C, au contraire,
montre qu'une validation explicite peut supprimer les injections executees dans
le scenario mesure.

Ce resultat est utile au-dela de la simulation. Tout agent d'entreprise connecte
a des outils doit traiter les sources externes comme non fiables par defaut :
documents RAG, messages utilisateurs, evenements SSE, resultats d'API et traces
d'autres agents. La separation entre instructions systeme et donnees doit etre
maintenue par conception, pas seulement par consigne au modele.

\section{Limites des benchmarks}

\subsection{Taille des echantillons}

La campagne principale porte sur 10 tickets et les tests A/B sur trois paires.
Ces volumes suffisent pour valider le pipeline et observer des tendances, mais
pas pour conclure statistiquement sur la performance generale. Le test de biais
de position illustre bien cette limite : la p-value de 0,5034 ne montre pas de
biais, mais l'echantillon reste trop petit pour garantir son absence.

\subsection{Dependance aux juges LLM}

Le juge \code{gpt-4o} apporte une evaluation scalable, mais il reste un modele.
Il peut avoir ses propres biais, variations et preferences de style. Le projet
reduit ce risque avec des tests de position et des criteres explicites, mais une
validation humaine reste necessaire pour une evaluation finale de qualite IT.

\subsection{Comparabilite des artefacts}

Tous les fichiers de resultats ne proviennent pas du meme protocole, de la meme
date ou des memes modeles. Par exemple, les scores historiques sur
\code{llama3.2:1b}, \code{gpt-3.5-turbo} et \code{mistral-7b} ne doivent pas
etre compares directement au benchmark \code{phi3:latest}. Ils servent de
repere, pas de preuve.

\subsection{Benchmark A2A non concluant}

Le fichier \file{a2a_agents_results.json} montre que l'orchestration A2A est en
place, mais chaque agent distant remonte une erreur. Le rapport ne peut donc pas
utiliser ces chiffres comme mesure de performance. C'est une limite importante
mais utile : elle indique que le pipeline doit mieux distinguer les erreurs de
transport, d'authentification, de protocole et d'evaluation.

\section{Qualite logicielle}

Le depot est bien testable : les 552 tests unitaires passent localement et les
LLM sont mocks dans les tests critiques. Cette decision est essentielle pour un
projet agentique. Sans mocks, la suite de tests serait lente, couteuse et
instable.

Quelques points restent a surveiller :

\begin{itemize}
  \item certaines donnees runtime et bases locales apparaissent modifiees dans
  l'arbre de travail ; il faut clarifier ce qui doit etre versionne ;
  \item la commande \code{bibops eval} est documentee, mais l'implementation
  locale semble s'appuyer sur un module historique \code{src.eval_bank} qui
  merite une verification avant livraison ;
  \item les artefacts de benchmark gagneraient a inclure systematiquement un
  identifiant de commit, la graine aleatoire, la version des prompts et la
  version exacte des modeles ;
  \item les graphiques sont deja presents, mais leur generation devrait etre
  incluse dans une commande unique de reproduction du rapport.
\end{itemize}

\section{Perspectives d'amelioration}

Les ameliorations prioritaires sont les suivantes.

\begin{enumerate}
  \item \textbf{Augmenter le volume d'evaluation.} Passer de 10 tickets a la
  totalite du jeu de 40 tickets, puis a un corpus plus large annote par des
  experts support.

  \item \textbf{Renforcer les prompts agentiques.} Imposer une reponse finale
  structuree, avec sources, niveau de confiance, actions et escalade.

  \item \textbf{Ameliorer le corpus RAG.} Ajouter des procedures internes plus
  completes, nettoyer les chunks, tracer les citations et evaluer la pertinence
  top-k.

  \item \textbf{Calibrer les scores avec des humains.} Comparer le juge LLM a un
  panel de correcteurs pour mesurer l'accord inter-juge.

  \item \textbf{Durcir la securite agentique.} Ajouter filtrage d'entree,
  sandbox d'outils, authentification des endpoints WeakProxy, separation claire
  entre donnees et instructions, et tests adversariaux automatises.

  \item \textbf{Industrialiser la reproductibilite.} Ajouter une commande de
  rapport qui regenere benchmarks, graphiques et PDF en consignant l'environnement
  d'execution.
\end{enumerate}

\section{Positionnement general}

\bibopsname{} est plus proche d'un banc d'industrialisation que d'une simple
demo LLM. C'est ce qui fait son interet. Le projet met en relation des themes
souvent traites separement : agents outilles, RAG, evaluation LLM, securite,
GreenOps, cout et simulation adversariale. Les resultats montrent que ces
dimensions peuvent etre mesurees ensemble, meme si la qualite finale n'est pas
encore au niveau requis.
